\documentclass[class=book, crop=false]{standalone}
% ==============================================================================
% SETUP.TEX - CONFIGURACIÓN GLOBAL DEL PROYECTO
% ==============================================================================
% Este archivo centraliza todos los paquetes y estilos.
% Debe incluirse en el preámbulo del libro principal y de los archivos standalone.
% \PassOptionsToPackage{gray}{xcolor}
% --- 1. CODIFICACIÓN E IDIOMA ---
% fix-cm habilita las fuentes Computer Modern en tamaños arbitrarios (por debajo
% de 5 pt, entre otros). Debe cargarse antes que fontenc.
\usepackage{fix-cm}
\usepackage[utf8]{inputenc}
\usepackage[T1]{fontenc}
\usepackage[spanish, es-tabla]{babel}  
\usepackage{csquotes} % Recomendado para babel y citas

\usepackage{import}
% mode=image: Usa el PDF pre-compilado, nunca intenta recompilar.
% Debes compilar cada figura manualmente antes de compilar el libro.
\usepackage[mode=image, subpreambles=false]{standalone}

% --- 2. MATEMÁTICAS Y CIENCIAS ---
\usepackage{amsmath, amssymb, amsfonts, mathtools}
\usepackage{enumitem}

% LaTeX trae \DeclareMathSizes{5}{5}{5}{5}, así que a 5 pt (\tiny) los subíndices
% salen del mismo tamaño que la base: en las etiquetas de figuras el M de
% $\omega_M$ queda desproporcionado. Se restablece la proporción habitual.
\DeclareMathSizes{5}{5}{3.7}{3}

% --- 3. DISEÑO Y GEOMETRÍA --- 
% Unificamos los márgenes para todo el documento
\usepackage[left=2.5cm,right=2.5cm,top=3cm,bottom=3cm]{geometry}
\makeatletter
\ifstandalone % Evita que geometry dañe el recorte de las figuras bajándolas 3cm
  \@ifclasswith{standalone}{crop=false}{
    % Es un capítulo/sección: Dejamos que geometry actúe.
  }{
    % Es una figura: Neutralizamos geometry para que no las recorte mal.
    \geometry{pass}
  }
\fi
\makeatother
\usepackage{parskip} % Espaciado entre párrafos sin sangría
\usepackage{float}   % Para usar [H] en figura

% Ajustes de flotantes para evitar espacios verticales exagerados
\renewcommand{\topfraction}{0.95}      % Permite que las figuras ocupen hasta el 95% superior
\renewcommand{\bottomfraction}{0.95}   % Permite que las figuras ocupen hasta el 95% inferior
\renewcommand{\textfraction}{0.05}     % Permite hojas con tan solo un 5% de texto
\renewcommand{\floatpagefraction}{0.8} % Requiere que una página de solo figuras esté 80% llena
\setcounter{topnumber}{4}
\setcounter{bottomnumber}{4}
\setcounter{totalnumber}{10}

% --- 4. GRÁFICOS Y TABLAS ---
\usepackage{graphicx}
\usepackage{booktabs} % Tablas profesionales
\usepackage{caption}  % Control de captions y \captionof
\usepackage{tikz}
\usetikzlibrary{arrows.meta, calc, angles, quotes, babel, backgrounds}
\usepackage{pgfplots}
\usepgfplotslibrary{groupplots}
\usepgfplotslibrary{external}
\usepgfplotslibrary{patchplots}
\pgfplotsset{compat=1.18}
\usepackage[american, straightvoltages]{circuitikz}

% --- 5. COLORES Y CAJAS ---

\usepackage[table]{xcolor}
\usepackage{tcolorbox}

% Definición de colores corporativos/estilo
\definecolor{utnblue}{RGB}{0, 51, 102}
\definecolor{codegreen}{rgb}{0,0.6,0}
\definecolor{codegray}{rgb}{0.5,0.5,0.5}
\definecolor{codepurple}{rgb}{0.58,0,0.82}
\definecolor{backcolour}{rgb}{0.96,0.96,0.96}

% --- 6. ENTORNOS PERSONALIZADOS (CAJAS) ---

% Caja "Resultado" (Azul Corporativo) — Definiciones, fórmulas clave, resultados importantes.
% Uso: \begin{resultado}[Fórmula de Euler] ... \end{resultado}
\newtcolorbox{resultado}[1][]{
    colback=utnblue!5!white,
    colframe=utnblue,
    title=\textbf{#1},
    fonttitle=\bfseries,
    sharp corners,
    boxrule=0.5mm
}

% Caja "Nota" (Gris) — Advertencias, aplicaciones, contexto secundario.
% Uso: \begin{nota}[Cuidado con la calculadora] ... \end{nota}
% Uso: \begin{nota} ... \end{nota}  → título por defecto "Nota"
\newtcolorbox{nota}[1][Nota]{
    colback=codegray!5!white,
    colframe=codegray,
    title=\textbf{#1},
    fonttitle=\bfseries,
    sharp corners,
    boxrule=0.5mm
}

% Caja inline para resaltar un valor y enlazarlo con su otra aparición en una tabla
% o en una cuenta: mismo color, mismo valor.
% Uso: \marcavalor{$4$}  o  \marcavalor[codegreen]{$4$}
\newtcbox{\marcavalor}[1][utnblue]{
    on line,
    boxsep=1pt, left=2pt, right=2pt, top=1pt, bottom=1pt,
    colback=#1!10!white,
    colframe=#1,
    boxrule=0.4pt,
    arc=2pt
}

% --- 7. CÓDIGO FUENTE (LISTINGS) ---
\usepackage{listings}

\lstdefinestyle{miestilo}{
    backgroundcolor=\color{backcolour},   
    commentstyle=\color{codegreen},
    keywordstyle=\color{magenta},
    numberstyle=\tiny\color{codegray},
    stringstyle=\color{codepurple},
    basicstyle=\ttfamily\footnotesize,
    breakatwhitespace=false,         
    breaklines=true,                 
    captionpos=b,                    
    keepspaces=true,                 
    numbers=left,                    
    numbersep=5pt,                  
    showspaces=false,                
    showstringspaces=false,
    showtabs=false,                  
    tabsize=2,
    frame=single,                    
    rulecolor=\color{codegray}
}

\lstset{style=miestilo}
% Soporte para tildes y ñ en bloques de código
\lstset{literate={ñ}{{\~n}}1 {á}{{\'a}}1 {é}{{\'e}}1 {í}{{\'i}}1 {ó}{{\'o}}1 {ú}{{\'u}}1}

% Operadores matemáticos personalizados

% Vector en sentido discreto (lista finita de coordenadas): negrita + flecha.
% Uso: \vect{v}, \vect{e}_k, \vect{0}.
\newcommand{\vect}[1]{\vec{\mathbf{#1}}}

% Fecha de versión y comando para capítulos standalone
\providecommand{\verdate}{\the\year.\ifnum\month<10 0\fi\the\month.\ifnum\day<10 0\fi\the\day}
\newcommand{\chapterversion}{%
    \vspace{-1.5cm}\begin{flushright}\small\textit{\color{codegray}Versión~\verdate}\end{flushright}\vspace{0.5cm}%
}

% --- 8. HIPERVÍNCULOS (DEBE SER EL ÚLTIMO PAQUETE) ---
\usepackage[colorlinks=true, linkcolor=utnblue, urlcolor=utnblue, citecolor=utnblue]{hyperref}

% --- 9. REFERENCIAS CRUZADAS ENTRE CAPÍTULOS STANDALONE ---
% Cuando se compila un capítulo standalone, xr-hyper lee los .aux de los
% otros capítulos (si existen) para resolver \ref{} a labels externos.
% En la compilación del libro completo este bloque no se activa.
\ifstandalone
  \usepackage{xr-hyper}
  \makeatletter
  \newcommand{\xrchap}[1]{%
    \edef\xrc@self{\detokenize{Capitulo#1}}%
    \edef\xrc@job{\jobname}%
    \ifx\xrc@self\xrc@job\else
      \IfFileExists{../#1capitulo/Capitulo#1.aux}%
        {\externaldocument{../#1capitulo/Capitulo#1}}{}%
    \fi
  }
  \makeatother
  \xrchap{01}\xrchap{02}\xrchap{03}\xrchap{04}
  \xrchap{05}\xrchap{06}\xrchap{07}\xrchap{08}
  \xrchap{09}\xrchap{10}\xrchap{11}\xrchap{12}
  \xrchap{13}
\fi

\begin{document}
\setcounter{chapter}{5}
\section{Serie trigonométrica}
\label{sec:serie_trigonometrica}

La serie exponencial~\eqref{eq:sf_sintesis_box} desarrollada
en la Sección~\ref{sec:serie_fourier} aplica a cualquier señal
periódica, real o compleja. Cuando $x(t)$ es real ---el caso de
prácticamente todas las señales que vamos a manipular en el
libro--- los coeficientes $c_k$ están emparejados por la simetría
conjugada $c_{-k} = \overline{c_k}$ probada en la
Subsección~\ref{subsec:prop_conjugacion}. Esa restricción permite
reescribir la serie en una forma equivalente con coeficientes
reales, la \textbf{serie trigonométrica}, que separa cada armónico
en una parte coseno y una parte seno. Aparece más a mano cuando la
señal tiene una simetría par o impar ---cada simetría anula
directamente una de las dos sumas--- y conecta de forma directa la
representación de Fourier con la descomposición de $x(t)$ en parte
par y parte impar.

% ==============================================================================
% 5.4.1 Conversión desde la exponencial
% ==============================================================================
\subsection{Conversión desde la exponencial}
\label{subsec:trig_conversion}

Partimos de la serie exponencial~\eqref{eq:sf_sintesis_box} y
aprovechamos la simetría conjugada $c_{-k} = \overline{c_k}$
de la Subsección~\ref{subsec:prop_conjugacion} para agrupar cada
par de términos con índices $k$ y $-k$. Los expandimos con la
identidad de Euler:
\[
    e^{jk\omega_0 t} = \cos(k\omega_0 t) + j\sin(k\omega_0 t),
    \qquad
    e^{-jk\omega_0 t} = \cos(k\omega_0 t) - j\sin(k\omega_0 t).
\]
Sumando los dos términos del par y agrupando coseno y seno por
separado,
\begin{equation*}
    c_k\,e^{jk\omega_0 t} + c_{-k}\,e^{-jk\omega_0 t}
    \;=\; (c_k + c_{-k})\cos(k\omega_0 t)
        + j\,(c_k - c_{-k})\sin(k\omega_0 t).
\end{equation*}
Aparecen entonces dos coeficientes, uno multiplicando al coseno y
otro al seno. Llamémoslos $a_k$ y $b_k$:
\begin{equation*}
    a_k = c_k + c_{-k},
    \qquad
    b_k = j\,(c_k - c_{-k}),
    \qquad k \geq 1.
\end{equation*}
Sumando sobre $k\geq 1$ y separando el término $k = 0$, la serie
queda
\begin{equation}
    x(t) \;=\; a_0 \;+\; \sum_{k=1}^{\infty}
        \bigl[a_k\cos(k\omega_0 t) + b_k\sin(k\omega_0 t)\bigr],
    \label{eq:sf_trig}
\end{equation}
con $a_0 = c_0$.

Para que la serie hable directamente de $x(t)$ ---sin pasar por los
$c_k$--- conviene tener fórmulas de $a_k$ y $b_k$ como integrales de
análisis. Reemplazando $c_k$ y $c_{-k}$ por sus
expresiones~\eqref{eq:ck_analisis_box} y sumando dentro de la
integral,
\begin{equation*}
    a_k \;=\; c_k + c_{-k}
        \;=\; \frac{1}{T_0}\int_{T_0} x(t)
            \bigl[e^{-jk\omega_0 t} + e^{jk\omega_0 t}\bigr]\,dt
        \;=\; \frac{2}{T_0}\int_{T_0} x(t)\cos(k\omega_0 t)\,dt,
\end{equation*}
donde la suma de las dos exponenciales vale $2\cos(k\omega_0 t)$ por
Euler. Para $b_k$ el cálculo es análogo:
\begin{equation*}
    b_k \;=\; j\,(c_k - c_{-k})
        \;=\; \frac{j}{T_0}\int_{T_0} x(t)
            \bigl[e^{-jk\omega_0 t} - e^{jk\omega_0 t}\bigr]\,dt
        \;=\; \frac{2}{T_0}\int_{T_0} x(t)\sin(k\omega_0 t)\,dt,
\end{equation*}
porque la diferencia entre paréntesis vale $-2j\sin(k\omega_0 t)$ y
el $j$ de afuera cancela el $-j$ de adentro. El término $a_0$ sale
de $a_0 = c_0$ y de la fórmula de análisis con $k = 0$:
\begin{equation*}
    a_0 \;=\; \frac{1}{T_0}\int_{T_0} x(t)\,dt,
\end{equation*}
es decir, el \textbf{valor medio} de $x$ sobre el período, en el
sentido del Capítulo~\ref{cap03}.

\paragraph{Convención del factor.}
La fórmula general de~$a_k$, particularizada en $k = 0$, daría el
doble del valor medio. Por convención se define $a_0$ con el factor
$1/T_0$ en lugar de $2/T_0$, de modo que el primer término de la
serie sea directamente el promedio.

\begin{resultado}[Serie trigonométrica]
Sea $x(t)$ una señal periódica real de período $T_0$ y frecuencia
angular $\omega_0 = 2\pi/T_0$. Su serie de Fourier admite la forma
\begin{equation}
    x(t) \;=\; a_0 \;+\; \sum_{k=1}^{\infty}
        \bigl[a_k\cos(k\omega_0 t) + b_k\sin(k\omega_0 t)\bigr],
    \label{eq:sf_trig_box}
\end{equation}
con coeficientes
\begin{align}
    a_0 &\;=\; \frac{1}{T_0}\int_{T_0} x(t)\,dt,
        \label{eq:trig_a0} \\[2pt]
    a_k &\;=\; \frac{2}{T_0}\int_{T_0} x(t)\cos(k\omega_0 t)\,dt,
        \qquad k \geq 1, \label{eq:trig_ak} \\[2pt]
    b_k &\;=\; \frac{2}{T_0}\int_{T_0} x(t)\sin(k\omega_0 t)\,dt,
        \qquad k \geq 1. \label{eq:trig_bk}
\end{align}
\end{resultado}

La relación con los coeficientes exponenciales queda como un
diccionario directo. Como $x$ es real, $c_{-k} = \overline{c_k}$, así
que $c_k + c_{-k} = 2\operatorname{Re}\{c_k\}$ y $c_k - c_{-k} =
2j\operatorname{Im}\{c_k\}$. Sustituyendo en las definiciones de
$a_k$ y $b_k$,
\begin{equation}
    a_k \;=\; 2\operatorname{Re}\{c_k\},
    \qquad
    b_k \;=\; -2\operatorname{Im}\{c_k\},
    \qquad k \geq 1.
    \label{eq:diccionario_trig_exp}
\end{equation}
A la inversa, $c_k = (a_k - jb_k)/2$ y $c_{-k} = \overline{c_k} =
(a_k + jb_k)/2$. Los $a_k$ leen la parte real del espectro complejo,
y los $b_k$ ---con un signo--- la parte imaginaria.

La figura~\ref{fig:base_trigonometrica} muestra las primeras tres
parejas de la base armónica trigonométrica
$\{\cos(k\omega_0 t),\,\sin(k\omega_0 t)\}$. Cualquier señal
periódica real, sumándole un valor medio $a_0$, se reconstruye como
combinación lineal de estas funciones (más sus armónicos
superiores), con pesos $a_k$ para los cosenos y $b_k$ para los
senos.

\begin{figure}[H]
    \centering
    \begin{minipage}{\linewidth}
        \centering
        \includestandalone[width=0.80\linewidth]{figures/fig_base_trigonometrica/fig_base_trigonometrica}
        \caption{Las primeras tres parejas de la base armónica
                 trigonométrica $\{\cos(k\omega_0 t),\,\sin(k\omega_0 t)\}$
                 sobre un período.
                 Fila superior, cosenos:
                 (a) $\cos(\omega_0 t)$, (b) $\cos(2\omega_0 t)$,
                 (c) $\cos(3\omega_0 t)$.
                 Fila inferior, senos:
                 (d) $\sin(\omega_0 t)$, (e) $\sin(2\omega_0 t)$,
                 (f) $\sin(3\omega_0 t)$. La serie trigonométrica
                 reescribe cualquier señal periódica real como
                 combinación lineal de estas funciones, sus armónicos
                 superiores, y la constante $a_0$.}
        \label{fig:base_trigonometrica}
    \end{minipage}
\end{figure}

% ==============================================================================
% 5.4.2 Lectura por simetría
% ==============================================================================
\subsection{Lectura por simetría}
\label{subsec:trig_simetrias}

La Subsección~\ref{subsec:par_impar} estableció que toda señal
admite una descomposición única como suma de una parte par y una
parte impar. La forma trigonométrica~\eqref{eq:sf_trig_box} se conecta
directamente con esa descomposición: los cosenos $\cos(k\omega_0 t)$
son funciones pares y los senos $\sin(k\omega_0 t)$ son impares, así
que la primera suma es una señal par y la segunda es una señal
impar. La paridad de $x$ se traduce entonces de forma inmediata en
qué coeficientes pueden ser distintos de cero.

Si $x(t)$ es \textbf{par}, todos los $b_k$ se anulan. La verificación
es inmediata desde~\eqref{eq:trig_bk}: el integrando es
$x(t)\sin(k\omega_0 t)$, producto de una función par por una impar,
que es impar (regla de productos de la
Subsección~\ref{subsubsec:productos_paridad}).
Integrar una función impar sobre un intervalo simétrico da
cero~\eqref{eq:integral_impar}, así que la integral
de~\eqref{eq:trig_bk} se anula para todo $k\geq 1$.

Si $x(t)$ es \textbf{impar}, los que se anulan son $a_0$ y los
$a_k$. El argumento es paralelo: en~\eqref{eq:trig_a0} el integrando
es $x(t)$, una función impar, y en~\eqref{eq:trig_ak} el integrando
es $x(t)\cos(k\omega_0 t)$, también impar por la misma regla de
productos. Las dos integrales se anulan sobre el intervalo simétrico.

\begin{resultado}[Simetrías par e impar]
Sea $x(t)$ una señal periódica real, con coeficientes
trigonométricos $a_0$, $a_k$, $b_k$.
\begin{itemize}
    \item Si $x$ es \textbf{par}, entonces $b_k = 0$ para todo
          $k \geq 1$ y la serie se reduce a una suma de cosenos:
          \[
              x(t) \;=\; a_0 \;+\; \sum_{k=1}^{\infty}
                  a_k\cos(k\omega_0 t).
          \]
    \item Si $x$ es \textbf{impar}, entonces $a_0 = 0$ y $a_k = 0$
          para todo $k \geq 1$, y la serie se reduce a una suma de
          senos:
          \[
              x(t) \;=\; \sum_{k=1}^{\infty}
                  b_k\sin(k\omega_0 t).
          \]
\end{itemize}
\end{resultado}

Para una señal que no es ni par ni impar, la
descomposición~\eqref{eq:descomposicion_par_impar} sigue operando: los
$a_k$ son los coeficientes trigonométricos de la parte par de $x$, y
los $b_k$ los de la parte impar. La forma trigonométrica es esa
descomposición leída sobre la base armónica.

% ==============================================================================
% 5.4.3 Ejemplos
% ==============================================================================
\subsection{Ejemplos}
\label{subsec:ejemplos_trig}

Aplicamos la receta trigonométrica a los dos ejemplos que ya
trabajamos en exponencial en la Sección~\ref{sec:serie_fourier}: el
tren de pulsos rectangulares y el diente de sierra. Para cada uno
calculamos los coeficientes $a_0$, $a_k$, $b_k$ directamente desde
las fórmulas~\eqref{eq:trig_a0}--\eqref{eq:trig_bk} y al final
verificamos que coinciden con los $c_k$ ya conocidos por el
diccionario~\eqref{eq:diccionario_trig_exp}.

\paragraph{Pulso rectangular.}
Retomamos el tren de pulsos de la
Subsección~\ref{subsec:ejemplo_pulso}: $x(t) = 1$ sobre la ventana
centrada $|t|\leq 1/2$ y $0$ en el resto del período $T_0 = 4$, con
frecuencia angular $\omega_0 = \pi/2$.

Como $x$ es real y par, la regla de simetría de la
Subsección~\ref{subsec:trig_simetrias} da inmediatamente
$b_k = 0$ para todo $k\geq 1$. El coeficiente $a_0$ es el valor
medio de $x$ sobre el período. Por inspección ---ventana de ancho
$1$ y amplitud $1$, dividida por el período de longitud $4$---,
\begin{equation*}
    a_0 \;=\; \frac{1}{4}.
\end{equation*}

Para los $a_k$ con $k\geq 1$, la fórmula~\eqref{eq:trig_ak} se
reduce a la ventana central porque $x = 0$ fuera de ella:
\begin{equation*}
    a_k \;=\; \frac{2}{4}\int_{-1/2}^{1/2}\cos(k\omega_0 t)\,dt
        \;=\; \frac{1}{2}\,
            \left[\frac{\sin(k\omega_0 t)}{k\omega_0}\right]_{-1/2}^{1/2}
        \;=\; \frac{1}{2}\cdot\frac{2\sin(k\omega_0/2)}{k\omega_0}
        \;=\; \frac{\sin(k\omega_0/2)}{k\omega_0},
\end{equation*}
donde usamos $\sin(-\theta) = -\sin(\theta)$ para simplificar la
diferencia en el numerador. Para reconocer el seno cardinal
$\operatorname{sinc}(x) = \sin(x)/x$ ---introducido en la
Subsección~\ref{subsec:sinc}--- basta notar que el denominador es el
doble del argumento del seno:
\begin{equation*}
    a_k \;=\; \frac{1}{2}\,\frac{\sin(k\omega_0/2)}{k\omega_0/2}
        \;=\; \frac{1}{2}\,\operatorname{sinc}\!\Bigl(\frac{k\omega_0}{2}\Bigr).
\end{equation*}
Sustituyendo $\omega_0 = \pi/2$,
\begin{equation}
    a_k \;=\; \frac{1}{2}\,\operatorname{sinc}\!\Bigl(\frac{k\pi}{4}\Bigr),
    \qquad k\geq 1.
    \label{eq:ak_pulso}
\end{equation}
La serie trigonométrica del pulso queda
\begin{equation*}
    x(t) \;=\; \frac{1}{4} \;+\;
        \sum_{k=1}^{\infty}
        \frac{1}{2}\,\operatorname{sinc}\!\Bigl(\frac{k\pi}{4}\Bigr)
        \cos\!\Bigl(\frac{k\pi t}{2}\Bigr).
\end{equation*}

\textbf{Verificación.} El diccionario~\eqref{eq:diccionario_trig_exp}
predice $a_k = 2\operatorname{Re}\{c_k\}$ y
$b_k = -2\operatorname{Im}\{c_k\}$. Los $c_k$ del pulso calculados
en~\eqref{eq:ck_pulso} son $c_k = (1/4)\operatorname{sinc}(k\pi/4)$,
reales para todo $k$. Entonces
$a_k = 2\cdot(1/4)\operatorname{sinc}(k\pi/4) = (1/2)\operatorname{sinc}(k\pi/4)$
y $b_k = 0$, en concordancia con los valores obtenidos por
integración directa.

\begin{figure}[H]
    \centering
    \begin{minipage}{\linewidth}
        \centering
        \includestandalone[width=0.70\linewidth]{figures/fig_aprox_trig_pulso/fig_aprox_trig_pulso}
        \caption{Aproximación trigonométrica del tren de pulsos
                 rectangulares. Cada fila corresponde a un orden
                 armónico $k = 1, 2, \dots, 5$. Columna izquierda:
                 armónico individual $a_k\cos(k\omega_0 t)$ (rojo);
                 en $k=4$ el armónico se anula porque
                 $\operatorname{sinc}(\pi) = 0$. Columna derecha:
                 suma parcial $\hat x_k(t) = a_0 + \sum_{m=1}^{k}
                 a_m\cos(m\omega_0 t)$ (rojo) superpuesta a la señal
                 original (azul). (a)(b)~$k=1$; (c)(d)~$k=2$;
                 (e)(f)~$k=3$; (g)(h)~$k=4$; (i)(j)~$k=5$.}
        \label{fig:aprox_trig_pulso}
    \end{minipage}
\end{figure}

\paragraph{Diente de sierra.}
Retomamos el diente de la
Subsección~\ref{subsec:ejemplo_sierra}: $x(t) = t/2$ sobre el
intervalo $[0, 4)$, periódico con $T_0 = 4$ y $\omega_0 = \pi/2$. A
diferencia del pulso, esta señal no es par ni impar tal cual está
definida, así que calculamos los tres coeficientes uno por uno.

El valor medio sale directo, ya sea por la lectura geométrica de la
rampa ---promedio entre el inicio $0$ y el final $2$--- o
aplicando~\eqref{eq:trig_a0}:
\begin{equation*}
    a_0 \;=\; \frac{1}{4}\int_0^4 \frac{t}{2}\,dt
        \;=\; \frac{1}{8}\cdot\frac{t^2}{2}\bigg|_0^4
        \;=\; 1.
\end{equation*}

Los $a_k$ y $b_k$ con $k\geq 1$ piden integración por partes. Para
los $a_k$, con $u = t$ y $v = \sin(k\omega_0 t)/(k\omega_0)$,
\begin{equation*}
    a_k \;=\; \frac{1}{4}\int_0^4 t\cos(k\omega_0 t)\,dt
        \;=\; \frac{1}{4}\Biggl\{
            \left[\frac{t\sin(k\omega_0 t)}{k\omega_0}\right]_0^4
            - \int_0^4 \frac{\sin(k\omega_0 t)}{k\omega_0}\,dt
            \Biggr\}.
\end{equation*}
El término de borde se anula porque
$\sin(k\omega_0\cdot 4) = \sin(2k\pi) = 0$ para todo $k$ entero, y
la integral restante es la de un seno armónico sobre un período
completo, que vale cero. Entonces $a_k = 0$ para todo $k\geq 1$.

Para los $b_k$, el procedimiento es análogo, con $u = t$ y
$v = -\cos(k\omega_0 t)/(k\omega_0)$:
\begin{equation*}
    b_k \;=\; \frac{1}{4}\int_0^4 t\sin(k\omega_0 t)\,dt
        \;=\; \frac{1}{4}\Biggl\{
            \left[-\frac{t\cos(k\omega_0 t)}{k\omega_0}\right]_0^4
            + \int_0^4 \frac{\cos(k\omega_0 t)}{k\omega_0}\,dt
            \Biggr\}.
\end{equation*}
El término de borde evaluado entre $0$ y $4$ vale
$-4\cos(2k\pi)/(k\omega_0) = -4/(k\omega_0)$, y la integral restante
es la de un coseno armónico sobre un período completo, que vale
cero. Reuniendo y reemplazando $\omega_0 = \pi/2$,
\begin{equation}
    b_k \;=\; \frac{1}{4}\cdot\frac{-4}{k\omega_0}
        \;=\; -\frac{2}{k\pi}, \qquad k\geq 1.
    \label{eq:bk_sierra}
\end{equation}
La serie trigonométrica del diente queda
\begin{equation*}
    x(t) \;=\; 1 \;-\; \frac{2}{\pi}\sum_{k=1}^{\infty}
        \frac{1}{k}\sin\!\Bigl(\frac{k\pi t}{2}\Bigr),
\end{equation*}
la serie clásica del diente de sierra.

\textbf{Verificación.} Los $c_k$ del diente calculados
en~\eqref{eq:ck_sierra} son $c_0 = 1$ y $c_k = j/(\pi k)$ para
$k\neq 0$, puramente imaginarios. El
diccionario~\eqref{eq:diccionario_trig_exp} predice
$a_k = 2\operatorname{Re}\{c_k\} = 0$ para $k\geq 1$, y
$b_k = -2\operatorname{Im}\{c_k\} = -2/(\pi k)$, en concordancia con
los valores obtenidos por integración directa.

\begin{figure}[H]
    \centering
    \begin{minipage}{\linewidth}
        \centering
        \includestandalone[width=0.70\linewidth]{figures/fig_aprox_trig_sierra/fig_aprox_trig_sierra}
        \caption{Aproximación trigonométrica del diente de sierra.
                 Cada fila corresponde a un orden armónico
                 $k = 1, 2, \dots, 5$. Columna izquierda: armónico
                 individual $b_k\sin(k\omega_0 t)$ (rojo). Columna
                 derecha: suma parcial $\hat x_k(t) = a_0 +
                 \sum_{m=1}^{k} b_m\sin(m\omega_0 t)$ (rojo)
                 superpuesta a la señal original (azul).
                 (a)(b)~$k=1$; (c)(d)~$k=2$; (e)(f)~$k=3$;
                 (g)(h)~$k=4$; (i)(j)~$k=5$.}
        \label{fig:aprox_trig_sierra}
    \end{minipage}
\end{figure}

\end{document}
